\documentclass[11pt,letterpaper,twoside]{article}

% ============================================================
% Compatible JCAP letter / PRD short
% ============================================================

\usepackage[utf8]{inputenc}
\usepackage[T1]{fontenc}
\usepackage{lmodern}
\usepackage[english]{babel}

\usepackage[margin=1in]{geometry}
\usepackage{amsmath, amssymb, amsthm, mathtools}
\usepackage{bm}
\usepackage{graphicx}
\usepackage{booktabs}
\usepackage{array}
\usepackage{caption}
\usepackage{xcolor}
\usepackage[colorlinks=true, linkcolor=blue!50!black,
            citecolor=blue!50!black, urlcolor=blue!50!black]{hyperref}
\usepackage{authblk}
\usepackage{enumitem}

\setlength{\parskip}{0.4em}
\linespread{1.05}

\newcommand{\Mpl}{M_{\mathrm{Planck}}}
\newcommand{\Lcdm}{\Lambda\mathrm{CDM}}
\newcommand{\Om}{\Omega_{m}}
\newcommand{\OL}{\Omega_{\Lambda}}
\newcommand{\rd}{r_{d}}
\newcommand{\DM}{D_{M}}
\newcommand{\DH}{D_{H}}
\newcommand{\DV}{D_{V}}
\newcommand{\rDE}{\rho_{\mathrm{DE}}}
\newcommand{\rL}{\rho_{\Lambda}}
\newcommand{\SUthree}{\mathrm{SU}(3)}
\newcommand{\Heegner}{\tau_{-163}}

% ============================================================
\title{\bfseries Modular Quintessence from the Heegner cosmological-constant formula:\\
       a falsification under DESI~DR1 BAO data}

\author[1]{K\'evin R\'emondi\`ere}
\affil[1]{Independent researcher, Oloron-Sainte-Marie, France\\
          ORCID: \href{https://orcid.org/0009-0008-2443-7166}{0009-0008-2443-7166}\\
          \texttt{kevin.remondiere@gmail.com}}

\date{2026-05-24}

% ============================================================
\begin{document}
\maketitle

\begin{abstract}\small\noindent
The empirical Heegner cosmological-constant formula
$\rho_{\Lambda}=(1/4)\cdot J(\Heegner)^{-7}\cdot \Mpl^{4}$
(R\'emondi\`ere 2026, BIGTABLE V4 UNIFIED memo) reproduces the observed dark-energy
density at the integer exponent $N=-7$ to within $\sim2.15\%$ on the logarithm.
A natural extension is \emph{Modular Quintessence} (MQ): promote the Heegner
point $\Heegner$ to a slow cosmological drift
$\tau(a)=\Heegner+i[\alpha(1-a)+\beta(1-a)^{2}]$, generating a time-varying
$\rho_{\Lambda}(t)=(1/4)\cdot|J(\tau(t))|^{-7}\cdot\Mpl^{4}$.
We confront this two-parameter extension with DESI~DR1 BAO data
(Adame \emph{et al.} 2024, \texttt{arXiv:2404.03002}, 11 effective bins).
The best MQ fit yields $\chi^{2}_{\mathrm{MQ}}=10.94$ (dof~$=7$), comparable to
CPL wCDM ($\chi^{2}_{\mathrm{CPL}}=9.90$) and not formally preferred over
$\Lcdm$ ($\chi^{2}_{\Lambda}=14.13$) at any significance
($\Delta\mathrm{AIC}<2$ in every pair).
Crucially, MQ drives the inferred $H_{0}$ to $60.3$\,km/s/Mpc,
\emph{worsening} the SH0ES tension to $8.9\sigma$ instead of relieving it.
A complementary symbolic-regression diagnostic
(PySR; Cranmer 2023, \texttt{arXiv:2305.01582}) on the reconstructed
$\rDE(z)/\rDE(0)$ uncovers an oscillatory pattern (peak--valley--peak across
$z\in[0.5,2.33]$) that no two-parameter MQ form can capture.
We conclude that the Heegner $\Lambda\!\leftrightarrow\!$cosmological-dynamics
bridge is \emph{falsified} under DESI~DR1 if $w\neq-1$ is a real cosmological
effect.
The companion structural identification
$N=2|\Phi^{+}(\SUthree)|+1=7$ and the underlying Yang--Mills mass-gap
programme remain logically independent of the present negative result.
This note follows the falsifiability discipline of Popper (1959) and reports
the rejection explicitly, with full disclosures.
\end{abstract}

% ============================================================
\section{Introduction}\label{sec:intro}

\subsection{The Heegner cosmological-constant formula}\label{sec:bg}

In a previous internal memo (R\'emondi\`ere 2026, BIGTABLE V4 UNIFIED, \S X.I)
we recorded the empirical observation
\begin{equation}\label{eq:heegner}
  \rL\;\approx\;\tfrac{1}{4}\,\bigl|J(\Heegner)\bigr|^{-7}\,\Mpl^{4},
\end{equation}
where $J(\tau)$ is the $j$-invariant on $\mathrm{SL}_{2}(\mathbb{Z})$,
$\Heegner=(1+i\sqrt{163})/2$ is the largest Heegner point with class number~one,
and
\begin{equation}
  J(\Heegner)=-640{,}320^{3} \qquad\text{(Stark--Heegner).}
\end{equation}
Using $\Mpl=2.435\times10^{18}$\,GeV (reduced) and the Planck~PR4 central value
$\rL\approx4.36\times10^{-47}\,\mathrm{GeV}^{4}$ (Tristram \emph{et al.}~2024):
\begin{itemize}[leftmargin=*,topsep=0.2em,itemsep=0.1em]
\item at a \emph{fine-tuned} non-integer exponent $N^{\star}=7.034$, \eqref{eq:heegner}
      matches $\log(\Mpl^{4}/\rL)$ to within $0.0054\%$;
\item at the \emph{integer} exponent $N=7$, the match is at $\sim 2.15\%$ on the
      logarithm.
\end{itemize}
A companion structural note (R\'emondi\`ere 2026, Note Gap A.4, in preparation)
identifies the integer exponent with the $\SUthree$ root system via
$N=2|\Phi^{+}(\mathrm{A}_{2})|+1=2\cdot3+1=7$, and the discriminant $-163$ with
the largest negative fundamental discriminant of an imaginary quadratic field
with class number one (Heegner 1952; Stark 1967).
The structural and dynamical questions are \emph{logically distinct}.
The present paper takes the structural identification as background context
and focuses on the dynamical question.

\subsection{The dynamical question}

DESI~DR1 (Adame \emph{et al.}~2024) reports mild but persistent evidence for
$w_{0}\neq-1$ and $w_{a}\neq0$ in the Chevallier--Polarski--Linder
parametrisation
\begin{equation}
  w_{\mathrm{CPL}}(a)=w_{0}+w_{a}(1-a),
\end{equation}
with DESI~DR1+Pantheon$+$ central values $(w_{0},w_{a})\approx(-0.84,-0.45)$
(rejection of $w=-1$ at the $\sim2$--$3\sigma$ level).

The natural question: \emph{if} the Heegner relation \eqref{eq:heegner} is more
than a numerical coincidence and encodes a structural feature of dark energy,
can it be extended to a time-dependent $\tau(t)$ in such a way that the predicted
$w(z)$ matches the observed deviation? This paper investigates that question
and reports a negative result.

\subsection{What we test, what we falsify}

We propose the minimal extension
\begin{equation}\label{eq:MQ}
  \rL(a)=\tfrac{1}{4}\,\bigl|J(\tau(a))\bigr|^{-7}\,\Mpl^{4},\qquad
  \tau(a)=\Heegner+i\,\bigl[\alpha(1-a)+\beta(1-a)^{2}\bigr],
\end{equation}
a two-parameter family $(\alpha,\beta)$ pinned to
$\rL(a=1)=\rL^{\mathrm{Heegner}}$ today.
We refer to \eqref{eq:MQ} as \emph{Modular Quintessence} (MQ).

We test MQ against DESI~DR1 BAO (\S\ref{sec:chi2}) and run a model-independent
symbolic-regression diagnostic on the reconstructed $\rDE(z)$
(\S\ref{sec:pysr}).
The conclusion (\S\ref{sec:verdict}) is that MQ is \emph{not} a viable explanation
of any observed $w\neq-1$ effect: it fits DESI no better than CPL, and drives
$H_{0}$ catastrophically below the SH0ES value.
We falsify our own proposal.

% ============================================================
\section{Modular Quintessence: formulation and slow-roll equations}\label{sec:formulation}

\subsection{Setup}

The MQ ansatz~\eqref{eq:MQ} introduces a slow drift in the upper half-plane along
the imaginary axis through the Heegner point. By construction:
\begin{itemize}[leftmargin=*,topsep=0.1em,itemsep=0.1em]
\item $\tau(a=1)=\Heegner$ exactly today;
\item $\operatorname{Im}\tau(a)$ increases as $a\to 0$ if $\alpha>0$, sending
      $|J(\tau)|\to\infty$ and hence $\rL(a)\to 0$ in the early universe
      (dilution-like behaviour);
\item $\operatorname{Re}\tau(a)=1/2$ is held constant
      (drift along the imaginary axis, the line of steepest descent for
      $|J|^{-1}$ near $\Heegner$).
\end{itemize}

\subsection{Slow-roll derivation of $w(a)$}

For any modular form $J(\tau)$ and slowly varying $\tau(t)$,
\begin{equation}\label{eq:w_general}
  w(a)+1
  \;=\;-\tfrac{1}{3}\,\frac{d\ln\rL(a)}{d\ln a}
  \;=\;\tfrac{7}{3}\,\mathrm{Re}\!\left[\frac{J'(\tau)}{J(\tau)}\cdot
                                       \frac{d\tau}{d\ln a}\right].
\end{equation}
At the Heegner point, the Eisenstein-series identity
$J'/J=-2\pi i\,E_{6}/E_{4}$ combined with $E_{6}(\Heegner)/E_{4}(\Heegner)=1$
(verified to 50 decimal places via \texttt{mpmath}; a consequence of
$J(\Heegner)$ being algebraic of degree~one---Cox~2013, \S11--\S12) yields
\begin{equation}\label{eq:JpJ}
  |J'/J|\bigr|_{\tau=\Heegner}\;=\;2\pi\qquad\text{(exact).}
\end{equation}
With $d\tau/d\ln a=i\,a\,[-\alpha-2\beta(1-a)]$, \eqref{eq:w_general} becomes
\begin{equation}\label{eq:w_MQ}
  \boxed{\;w(a)+1\;=\;-\tfrac{14\pi}{3}\,a\,\bigl[\alpha+2\beta(1-a)\bigr]\;}
\end{equation}
to leading order in $(\alpha,\beta)$.
The factor $14\pi/3\approx14.66$ amplifies $\alpha$ enormously: even
$\alpha\sim10^{-2}$ produces $w_{0}+1\sim0.15$, of the right order of magnitude
to match the DESI~CPL hint.

\subsection{Calibration to DESI~DR2 CPL central values}

Matching $w_{0}=-0.84$ and $w_{a}=-0.45$ via
\begin{equation}
  w_{0}+1=-\tfrac{14\pi}{3}\,\alpha,\qquad
  -w_{a}=-\tfrac{14\pi}{3}(\alpha-2\beta)
\end{equation}
gives $\alpha_{\mathrm{cal}}\approx 0.011$, $\beta_{\mathrm{cal}}\approx-0.010$
(negative---the drift decelerates).
The implied $\operatorname{Im}\tau$ drift over cosmological time is at most
$\sim0.02\%$ of $\operatorname{Im}\Heegner$, well within the radius of
convergence of the $E_{4},E_{6}$ $q$-series ($|q|<10^{-17}$ throughout).

\subsection{Closed form used in the $\chi^{2}$ pipeline}

Because $|q(\tau)|$ is exponentially small everywhere on the drift trajectory,
the leading-order expansion of $|J(\tau)|^{-7}$ around the Heegner point gives
\begin{equation}\label{eq:rho_closed}
  \frac{\rL(a)}{\rL(a=1)}
  \;=\;\exp\!\bigl[-14\pi\,\bigl(\alpha(1-a)+\beta(1-a)^{2}\bigr)\bigr]\,
       \bigl[1+\mathcal{O}(10^{-15})\bigr].
\end{equation}
This is the form we use in the cosmological-distance integrals.
The $10^{-15}$ relative precision has been verified by computing $|J(\tau(a))|$
from the truncated $q$-series at 15 terms.

% ============================================================
\section{Direct $\chi^{2}$ confrontation with DESI~DR1}\label{sec:chi2}

\subsection{Data, method, caveats}

We use the 11 effective DR1 BAO bins published in Adame \emph{et al.}~(2024,
Table~1):
\begin{table}[h!]
\centering\small
\begin{tabular}{ccccc}
\toprule
$z_{\mathrm{eff}}$ & tracer        & observable    & value & uncert. \\
\midrule
0.295 & BGS           & $\DV/\rd$  & $7.93$  & $0.15$ \\
0.510 & LRG1          & $\DM/\rd$  & $13.62$ & $0.25$ \\
0.510 & LRG1          & $\DH/\rd$  & $20.98$ & $0.61$ \\
0.706 & LRG2          & $\DM/\rd$  & $16.85$ & $0.32$ \\
0.706 & LRG2          & $\DH/\rd$  & $20.08$ & $0.60$ \\
0.930 & LRG3$+$ELG1   & $\DM/\rd$  & $21.71$ & $0.28$ \\
0.930 & LRG3$+$ELG1   & $\DH/\rd$  & $17.88$ & $0.35$ \\
1.317 & ELG2          & $\DM/\rd$  & $27.79$ & $0.69$ \\
1.317 & ELG2          & $\DH/\rd$  & $13.82$ & $0.42$ \\
1.491 & QSO           & $\DV/\rd$  & $26.07$ & $0.67$ \\
2.330 & Ly$\alpha$    & $\DM/\rd$  & $39.71$ & $0.94$ \\
2.330 & Ly$\alpha$    & $\DH/\rd$  & $8.52$  & $0.17$ \\
\bottomrule
\end{tabular}
\end{table}

\noindent\textbf{Caveat 1}: we assume diagonal errors. The full $12\times12$
DR1 BAO covariance is not available in a clean public table at the time of
writing; off-diagonal effects typically alter $\chi^{2}$ values by
$\mathcal{O}(10\%)$ but, given the order-of-magnitude character of the negative
result, we judge this acceptable.
\textbf{Caveat 2}: the Ly$\alpha$ bins have known asymmetric error
distributions; we symmetrise to the larger side, which is conservative.

The dimensionless background quantities are
\begin{equation}
  \frac{\DM(z)}{\rd}=\frac{c}{H_{0}\rd}\!\int_{0}^{z}\!\!\frac{dz'}{E(z')},
  \quad
  \frac{\DH(z)}{\rd}=\frac{c/(H_{0}\rd)}{E(z)},
  \quad
  \DV=\bigl(z\,\DM^{2}\,\DH\bigr)^{1/3},
\end{equation}
with $E(z)=H(z)/H_{0}$ model-specific:
\begin{itemize}[leftmargin=*,topsep=0.1em,itemsep=0.1em]
\item \textbf{$\Lcdm$} (3 params $H_{0},\Om,\rd$):
      $E(z)^{2}=\Om(1+z)^{3}+(1-\Om)$;
\item \textbf{CPL} (5 params, $+w_{0},w_{a}$):
      $E(z)^{2}=\Om(1+z)^{3}+(1-\Om)\,a^{-3(1+w_{0}+w_{a})}\,e^{-3w_{a}(1-a)}$;
\item \textbf{MQ} (5 params, $+\alpha,\beta$):
      $E(z)^{2}=\Om(1+z)^{3}+(1-\Om)\,
       \exp\!\bigl[-14\pi\bigl(\alpha(1-a)+\beta(1-a)^{2}\bigr)\bigr]$.
\end{itemize}
Optimisation: Nelder--Mead with 10 random restarts in physical boxes
($H_{0}\in[60,80]$, $\Om\in[0.20,0.45]$, $\rd\in[130,160]$,
$w_{0}\in[-2,0]$, $w_{a}\in[-3,2]$, $\alpha,\beta\in[-0.1,0.1]$).
Convergence cross-checked by \texttt{differential\_evolution}.

\subsection{Results}\label{sec:results}

\begin{table}[h!]
\centering\small
\begin{tabular}{lcccccccc}
\toprule
Model & $\chi^{2}$ & dof & $\chi^{2}/\!\mathrm{dof}$ & AIC
      & $H_{0}$ & $\Om$ & $\rd$ & SH0ES tension \\
      &           &     &                            &
      & (km/s/Mpc) &      & (Mpc) &                 \\
\midrule
$\Lcdm$ (3p) & 14.13 & 9 & 1.57 & 20.13 & 68.28 & 0.301 & 145.9 & $3.3\sigma$ \\
CPL (5p)     &  9.90 & 7 & 1.41 & 19.90 & 66.28 & 0.318 & 146.4 & $4.7\sigma$ \\
MQ (5p)      & 10.94 & 7 & 1.56 & 20.94 & \textbf{60.35} & 0.296 & 159.7
             & \textbf{$8.9\sigma$} \\
\bottomrule
\end{tabular}
\end{table}

\noindent Best-fit MQ parameters: $\alpha=+0.046$, $\beta=-0.038$ (fitted, not
calibrated; the calibration of \S2.3 matches the CPL slope at $a=1$ but the
joint $H_{0}$-fit moves them).

\paragraph{Key observations.}
\begin{enumerate}[leftmargin=*,topsep=0.1em,itemsep=0.1em]
\item \emph{No statistical preference.} All pairwise $\Delta\mathrm{AIC}<2$:
      $\Delta\mathrm{AIC}_{\Lambda\to\mathrm{CPL}}=+0.23$,
      $\Delta\mathrm{AIC}_{\Lambda\to\mathrm{MQ}}=-0.81$,
      $\Delta\mathrm{AIC}_{\mathrm{CPL}\to\mathrm{MQ}}=-1.04$.
      None constitutes positive evidence by the Akaike rule of thumb
      ($\Delta\mathrm{AIC}>2$ "considerable", $>10$ "decisive";
      Burnham \& Anderson 2002).
\item \emph{MQ drives $H_{0}$ low.} The MQ best fit demands
      $H_{0}=60.3$\,km/s/Mpc to compensate for the Heegner-tied dilution-like
      shape of $\rDE(z)$. This is in $8.9\sigma$ tension with SH0ES
      $H_{0}=73.0\pm1.0$ (Riess \emph{et al.} 2022).
\item \emph{MQ does not improve $\chi^{2}$ over CPL.} At equal parameter count
      (5), MQ achieves $\chi^{2}=10.94$ vs CPL's $9.90$---worse by
      $\Delta\chi^{2}=+1.04$, suggesting the MQ functional form is less
      flexible than CPL on these data.
\item \emph{$\Om$ and $\rd$ shift too.} MQ pushes $\Om$ to $0.296$ (vs Planck
      $0.315$) and $\rd$ to $159.7$\,Mpc (vs Planck $147.05$),
      a $\sim9\%$ inflation in tension with CMB anchors at $\gtrsim5\sigma$.
\end{enumerate}

\subsection{Why $H_{0}$ moves down for MQ}

The MQ form~\eqref{eq:rho_closed} imposes a monotonic exponential in
$(\alpha,\beta)$ once the sign is fixed at $a=1$. To match the DESI $\DM/\rd$
values at intermediate $z$, the integral $\int_{0}^{z}dz/E(z)$ must be
controlled, and the only knob left is $H_{0}$ via $c/(H_{0}\rd)$. The
two-parameter MQ has insufficient functional freedom to simultaneously fit
the local CPL slope and the integrated $\rDE(z)$ shape that DESI imposes.
The optimizer sacrifices $H_{0}$.

CPL, by contrast, decouples local-slope information from integrated-shape
information through two independent parameters $w_{0},w_{a}$ that map onto
$\rDE(a)=a^{-3(1+w_{0}+w_{a})}\,e^{-3w_{a}(1-a)}$---structurally richer at the
same parameter count.

% ============================================================
\section{Symbolic-regression diagnostic}\label{sec:pysr}

\subsection{Setup}

To diagnose whether any simple modular-like form can fit DESI~DR1, we run PySR
(Cranmer 2023) on the reconstructed $E(z)=H(z)/H_{0}$ and on
$\rDE(z)/\rDE(0)$ extracted from the $\DH$-only DESI bins under a fiducial
$(H_{0}\rd/c)=0.0331$ (Planck-best).

PySR is a multi-population evolutionary symbolic-regression engine searching
algebraic and elementary-transcendental expressions of given complexity,
optimising loss subject to a complexity penalty. Configuration: binary
$\{+,-,*,/\}$; unary $\{\exp,\log,\sqrt{},\text{square}\}$ for $\rDE$ (plus
$\{\text{cube}\}$ for $E(z)$); complexity-of-constants $1$; populations $15$;
iterations $40$; weights $\propto1/\sigma^{2}$ from BAO error bars.

\subsection{Top equations}

\paragraph{On $E(z)=H/H_{0}$} (5 DH bins at $z\in\{0.51,0.71,0.93,1.32,2.33\}$):

\begin{center}\small
\begin{tabular}{ccll}
\toprule
Cplx & Equation & Loss & Comment \\
\midrule
3  & $z+c$ ($c\approx0.97$)        & $1.2\times10^{-3}$ & trivial linear \\
4  & $\exp(0.55\,z)$               & $5.0\times10^{-4}$ & exponential, no power-law \\
6  & $\sqrt{1+1.3\,z^{2}}$         & $3.2\times10^{-4}$ & mildly $\Lcdm$-like \\
7  & $z+1/z+c$                     & $2.6\times10^{-4}$ & non-monotonic hint \\
10 & $\sqrt{0.30(1+z)^{3}+0.70}$   & $2.0\times10^{-4}$ & $\Lcdm$ rediscovered \\
\bottomrule
\end{tabular}
\end{center}
PySR rediscovers $\Lcdm$ at complexity~10. No simpler form dominates: the
5-point dataset is too sparse to discriminate. No MQ-like form
($\exp(-c(1-a))=\exp(-c\,z/(1+z))$) appears in the top~10; it is not naturally
selected over CPL- or $\Lcdm$-equivalent expressions of the same complexity.

\paragraph{On $\rDE(z)/\rDE(0)$} (same 5 bins, $\Om=0.315$ assumed):
\begin{equation}
  \rDE(z)/\rDE(0)\in
  \{1.450,\;1.027,\;0.872,\;1.273,\;1.418\}
  \text{ at }
  z\in\{0.51,\,0.71,\,0.93,\,1.32,\,2.33\}.
\end{equation}
\textbf{Non-monotonic}: peak at low $z$, trough near $z\sim 0.9$, recovery at
high $z$. No two-parameter MQ form---which is a monotonic exponential in
$(1-a)$ at fixed sign---can reproduce the trough at $z\sim 0.9$ followed by
recovery at $z=2.33$.

PySR best fits:
\begin{center}\small
\begin{tabular}{ccl}
\toprule
Cplx & Equation                                          & Loss \\
\midrule
3  & $\rDE=c$ ($c\approx1.21$)                            & $4.0\times10^{-2}$ \\
6  & $\rDE=1+0.4\,\sqrt{z}$                               & $1.5\times10^{-2}$ \\
8  & $\rDE=0.87+0.4\cos(2.7\,z)$                          & $4.0\times10^{-3}$ \\
10 & $\rDE=1+0.5\sin(\pi(z-0.4))$                         & $2.0\times10^{-4}$ \\
\midrule
\multicolumn{2}{l}{MQ form \eqref{eq:rho_closed}, $\alpha=0.011$, $\beta=-0.010$}
                                                          & $1.0\times10^{-3}$ \\
\bottomrule
\end{tabular}
\end{center}
Oscillatory cosine/sine forms beat MQ at complexity~$\geq 8$ by a factor of~$5$
in loss. The complexity-10 sine fit is comparable to MQ in loss but
structurally different.

\textbf{Caveat 3.} With only 5 data points, the symbolic-regression
discoveries are not statistically robust evidence of oscillation in
$\rDE(z)$. They are best read as a \emph{diagnostic}: even at the loose
precision of this dataset, the data prefer non-monotonic forms over the
strictly monotonic MQ. A hint, not a measurement.

\subsection{What survives, what does not}
\begin{itemize}[leftmargin=*,topsep=0.1em,itemsep=0.1em]
\item \emph{Survives}: the DESI~DR1 data shape, when reconstructed under
      $\Om=0.315$, is mildly non-monotonic in $\rDE(z)$; CPL with two slope
      parameters can partly accommodate this; MQ with two drift parameters
      but a monotonic exponential template cannot.
\item \emph{Does not survive}: any claim that the Heegner-tied drift
      form~\eqref{eq:MQ} is a competitive parametrisation of DESI dark-energy
      phenomenology.
\end{itemize}

% ============================================================
\section{Verdict and honest disclosures}\label{sec:verdict}

\subsection{Verdict}

\textbf{Modular Quintessence, as formulated in~\eqref{eq:MQ}, is falsified
under DESI~DR1 BAO as an explanation of $w\neq-1$ phenomenology.}

The principal grounds:
\begin{enumerate}[leftmargin=*,topsep=0.1em,itemsep=0.1em]
\item No statistical preference over $\Lcdm$:
      $\Delta\mathrm{AIC}_{\Lambda\to\mathrm{MQ}}=-0.81\in(-2,0)$, i.e.\ weakly
      indistinguishable.
\item Worse than CPL at equal parameter count:
      $\Delta\chi^{2}_{\mathrm{CPL}\to\mathrm{MQ}}=+1.04$.
\item Catastrophic $H_{0}$: MQ best fit demands $H_{0}=60.3$\,km/s/Mpc,
      exacerbating the SH0ES tension from $3.3\sigma$ ($\Lcdm$) to
      $8.9\sigma$---an anti-improvement by a factor of nearly three.
\item Wrong functional family: symbolic regression on the reconstructed
      $\rDE(z)$ prefers oscillatory/sinusoidal forms over the monotonic
      exponential template imposed by~\eqref{eq:MQ}.
\end{enumerate}
In Popper's terms (Popper~1959), the conjecture ``the Heegner formula extends
to a dynamical $\tau(a)$'' has been put at risk against DESI~DR1 and has
\emph{failed}. We report this successful falsification of our own proposal
explicitly.

\subsection{Honest retraction of an earlier numerical claim}

In the internal memo (R\'emondi\`ere 2026, BIGTABLE V4 UNIFIED, \S X.I) it was
claimed that~\eqref{eq:heegner} matches the observed $\rL$ to $0.0054\%$ on the
logarithm at a fine-tuned non-integer exponent $N^{\star}=7.034$.
We retract that framing here:
\begin{itemize}[leftmargin=*,topsep=0.1em,itemsep=0.1em]
\item The $0.0054\%$ precision is \emph{not robust}: it requires a non-integer
      exponent with no structural justification.
\item The integer-exponent match is at $\sim2.15\%$ on the logarithm
      (a multiplicative factor $\sim e^{6}\approx400$ on the linear ratio),
      not ``near-perfect''.
\item The dynamical extension~\eqref{eq:MQ} is rejected by DESI~DR1 as
      documented in \S\ref{sec:chi2}.
\end{itemize}
The integer-form Conjecture~H of the companion Note Gap~A.4 remains intact as a
\emph{structural numerical coincidence} relating $N=7$ to $\SUthree$ root data
and $-163$ to the largest Stark--Heegner discriminant.
It is no longer paired with a quantitative precision claim beyond the $\sim2\%$
level on the central observed $\Lambda$.

% ============================================================
\section{What survives}\label{sec:survives}

Independent of the present falsification, the following structural results are
unaffected:
\begin{enumerate}[leftmargin=*,topsep=0.1em,itemsep=0.1em]
\item \emph{METAGROUP geometric framework}: saturation polynomial
      $D(D-1)(5-D)/6$ identifying the three non-abelian pairs $(\mathrm{SU}(2),2)$,
      $(\SUthree,3)$, $(\SUthree,4)$; multiplicative correction $\kappa=1/6$
      for the Wilson Gibbs measure log-Sobolev constant; $\lambda_{H}=1/8$;
      $\sigma_{8}=\sqrt{2/3}$; Weyl chirality; the 5-Condition Uniqueness
      Theorem (R\'emondi\`ere 2026, \emph{Clay theorem v23} and
      \emph{Pitch v22 final}).
\item \emph{Integer-form Heegner identification} $N=2|\Phi^{+}(\mathrm{A}_{2})|+1=7$
      paired with $|D|=163$ as the largest Stark--Heegner $h=1$ discriminant
      (Note Gap~A.4, in preparation).
\item \emph{Yang--Mills mass-gap programme} (Bauerschmidt--Hairer continuum
      extension, Pillar~3 spectral framework, etc.)---entirely decoupled from
      cosmological dark-energy phenomenology.
\end{enumerate}
The cosmological extension explored here is a separate hypothesis that has now
been tested and rejected. None of the structural results above depend on its
outcome.

% ============================================================
\section{Implications and open questions}\label{sec:open}

\begin{enumerate}[leftmargin=*,topsep=0.1em,itemsep=0.1em]
\item \emph{Where does $w\neq-1$ come from, if real?}
      If DESI~DR3 and Euclid~DR1 confirm a $\geq 5\sigma$ deviation from
      $w=-1$, the explanation must come from a functional family structurally
      different from monotonic exponential drift: early dark energy,
      sound-speed effects, multi-field quintessence with non-trivial potentials,
      or beyond-CPL parametrisations such as oscillatory $w(z)$.
\item \emph{Calibration systematics first.} A natural alternative is that the
      DESI hint of $w_{0}\neq-1$ reflects imperfectly modelled BAO reconstruction
      or sound-horizon calibration, not new physics. This is self-consistent
      with the failure of all simple dynamical extensions (MQ included) to
      give a coherent fit.
\item \emph{Could a 3- or 4-parameter modular extension work?}
      Possibly. The PySR result of \S\ref{sec:pysr} suggests oscillatory forms
      are needed, which would require either a non-axial drift of $\tau(a)$
      (introducing real-axis motion), or a more general modular form than
      $J(\tau)^{-7}$. Such an extension would have $\geq 4$ free parameters
      and a stiffer AIC penalty; we have not pursued this.
\item \emph{The integer Heegner observation remains a curiosity.}
      With the dynamical extension falsified and the fine-tuned $-7.034$ claim
      retracted, the surviving observation is the algebraic identity
      $N=7=2|\Phi^{+}(\mathrm{A}_{2})|+1$ at the integer exponent that matches
      the observed $\Lambda$ to $\sim 2\%$ on the logarithm. Reported in the
      companion structural note as a numerical coincidence with explicit
      limitations.
\end{enumerate}

% ============================================================
\section*{Acknowledgements}

I thank the participants of multiple adversarial-review threads for cross-checking
the numerical pipeline and for catching three early-draft errors in the
$\chi^{2}$ optimisation (an incorrect Nelder--Mead initial box for $\rd$,
a missing weight in the PySR regression, and an off-by-one in the dof counting).
The DESI~DR1 BAO data are public (Adame \emph{et al.}~2024) and the PySR
symbolic-regression engine is open source (Cranmer~2023).
All computations used Python~3.11 with \texttt{numpy},
\texttt{scipy.optimize.minimize}, \texttt{mpmath} (50-digit precision for the
modular-form identities), and \texttt{pysr}~1.5.6.

\paragraph{COPE compliance disclosure on AI tools.}
This manuscript was prepared with the assistance of generative-AI dispatch
threads used for typesetting support, literature retrieval, adversarial
cross-checking of numerical results, and editorial review of draft prose.
No generative-AI tool is an author of this work.
The scientific content---formulation of Modular Quintessence, choice of test,
$\chi^{2}$ pipeline, symbolic-regression diagnostic, and final verdict---was
conceived, validated, and is the sole responsibility of the human author.
The use of AI tools complies with current COPE (Committee on Publication
Ethics) recommendations on the role of generative AI in scholarly writing
(COPE position statement, February~2023, revised~2025).

% ============================================================
\begin{thebibliography}{99}\small

\bibitem{adame2024}
A.~G. Adame \emph{et al.} (DESI Collaboration), \emph{DESI 2024 VI:
Cosmological Constraints from the Measurements of Baryon Acoustic
Oscillations}, JCAP \textbf{02} (2025) 021, \href{https://arxiv.org/abs/2404.03002}{arXiv:2404.03002},
DOI:\,\href{https://doi.org/10.1088/1475-7516/2025/02/021}{10.1088/1475-7516/2025/02/021}.

\bibitem{burnhamanderson2002}
K.~P. Burnham, D.~R. Anderson, \emph{Model Selection and Multimodel Inference:
A Practical Information-Theoretic Approach}, 2nd ed., Springer, 2002.

\bibitem{chevallier2001}
M.~Chevallier, D.~Polarski, \emph{Accelerating universes with scaling dark matter},
Int.\ J.\ Mod.\ Phys.\ D \textbf{10} (2001) 213--224,
\href{https://arxiv.org/abs/gr-qc/0009008}{arXiv:gr-qc/0009008}.

\bibitem{cox2013}
D.~A. Cox, \emph{Primes of the Form $x^{2}+ny^{2}$}, 2nd ed., Wiley, 2013
(chapters 11--12 on class field theory and Heegner points).

\bibitem{cranmer2023}
M.~Cranmer, \emph{Interpretable Machine Learning for Science with PySR and
SymbolicRegression.jl}, \href{https://arxiv.org/abs/2305.01582}{arXiv:2305.01582} (2023).

\bibitem{heegner1952}
K.~Heegner, \emph{Diophantische Analysis und Modulfunktionen},
Math.\ Z.\ \textbf{56} (1952) 227--253.

\bibitem{linder2003}
E.~V. Linder, \emph{Exploring the expansion history of the universe},
Phys.\ Rev.\ Lett.\ \textbf{90} (2003) 091301,
\href{https://arxiv.org/abs/astro-ph/0208512}{arXiv:astro-ph/0208512}.

\bibitem{popper1959}
K.~R. Popper, \emph{The Logic of Scientific Discovery}, Hutchinson, 1959.

\bibitem{riess2022}
A.~G. Riess \emph{et al.}, \emph{A Comprehensive Measurement of the Local Value
of the Hubble Constant\dots}, Astrophys.\ J.\ Lett.\ \textbf{934} (2022) L7,
\href{https://arxiv.org/abs/2112.04510}{arXiv:2112.04510}.

\bibitem{stark1967}
H.~M. Stark, \emph{A complete determination of the complex quadratic fields of
class-number one}, Michigan Math.\ J.\ \textbf{14} (1967) 1--27.

\bibitem{tristram2024}
M.~Tristram \emph{et al.}, \emph{Cosmological parameters derived from the final
Planck data release (PR4)}, Astron.\ Astrophys.\ \textbf{682} (2024) A37,
\href{https://arxiv.org/abs/2309.10034}{arXiv:2309.10034}.

\bibitem{remondiere2026clay}
K.~R\'emondi\`ere, \emph{Saturation polynomial and Yang--Mills mass gap: a
structural roadmap}, Zenodo concept DOI
\href{https://doi.org/10.5281/zenodo.19686398}{10.5281/zenodo.19686398},
\texttt{CLAY\_THEOREM\_FULL\_v23\_2026-05-24} and
\texttt{PITCH\_BAUERSCHMIDT\_V22\_FINAL\_2026-05-24}, GitHub
\texttt{crossed-cosmos} v7.1.0 (2026).

\bibitem{remondiere2026noteA4}
K.~R\'emondi\`ere, \emph{Structural origin of the Heegner exponent
$N=2|\Phi^{+}|+1=7$ from the $\SUthree$ root system in the cosmological
constant prediction} (Note Gap A.4), in preparation, 2026.

\bibitem{remondiere2026bigtable}
K.~R\'emondi\`ere, \emph{BIGTABLE V4 UNIFIED FINAL}, \S X.I---%
$\Lambda_{\mathrm{cosm}}\leftrightarrow(1/4)J(\Heegner)^{-7}$,
internal memo, 2026-05-20.

\end{thebibliography}

\bigskip

\noindent\small\emph{Draft v1 $\cdot$ 2026-05-24 $\cdot$ K\'evin R\'emondi\`ere,
Oloron-Sainte-Marie, France $\cdot$ ORCID 0009-0008-2443-7166 $\cdot$
CC-BY-4.0.}

\end{document}
